\section{The nonhereditary cluster-tilted $A_4$ certificate}
\label{sec:A4}

We return to the algebra which initiated the project.  Let
\begin{equation}
 Q:\quad
 1\xrightarrow a2\xrightarrow b4\xrightarrow c3\xrightarrow d2,
 \qquad
 A=\mathbb F_qQ/(bc,cd,db).
 \label{eq:A4-algebra}
\end{equation}
This is a nonhereditary cluster-tilted algebra of Dynkin type $A_4$.  Its
cluster category is modeled by a heptagon.

\subsection{The rigid atlas}

The heptagon has
\[
 \AIVDiagonals=14
\]
diagonals, one for each reduced indecomposable morphism object, and
\[
 \AIVTriangulations=42
\]
triangulations, one for each cluster/maximal rigid object.  An unordered
crossing pair is determined by four vertices, so there are
\[
 \AIVCrossingPairs=\binom74=35
\]
such pairs.  The Hall--skein ideal is generated by the two ordered relations
for each of these $35$ quadrilaterals, after the compatible Hall cocycle and
central boundary specialization have been fixed.

The previous seed computation used the initial exchange matrix
\begin{equation}
 B_0=
 \begin{pmatrix}
 0&-1&0&0\\
 1&0&1&-1\\
 0&-1&0&1\\
 0&1&-1&0
 \end{pmatrix},
 \qquad \det B_0=1.
 \label{eq:A4-B0}
\end{equation}
Thus the mutable matrix is full rank and admits the compatible form
\[
 \Lambda_0=(B_0^T)^{-1}.
\]
The global Hall cocycle of the preceding paper applies without frozen-rank
ambiguity.  The present quotient adds exactly the $35$ crossing relations
needed to close the algebra under incompatible products.

\subsection{The distinguished morphism object}

Consider
\begin{equation}
 C_{S_2}=(P_4\xrightarrow{\rho_b}P_2).
 \label{eq:A4-CS2}
\end{equation}
Its full homology is
\[
 H^{-1}(C_{S_2})\simeq S_3=\Omega_A^2S_2,
 \qquad
 H^0(C_{S_2})\simeq S_2,
\]
so the shifted-projective residue is zero and
\[
 I(C_{S_2})=e_2-e_4.
\]
In the natural morphism seed its cluster $g$-vector is
\[
 g(C_{S_2})=e_1-e_2+e_3.
\]
The corresponding classical cluster variable is
\begin{equation}
 CC_T^{\rm mor}(C_{S_2})
 =\frac{z_1z_3+z_4}{z_2},
 \label{eq:A4-classical-variable}
\end{equation}
and the based quantum expression is
\begin{equation}
 \qchar_{C_{S_2}}
 =X^{e_1-e_2+e_3}+X^{e_4-e_2}.
 \label{eq:A4-quantum-variable}
\end{equation}
The Hall--skein homomorphism sends the exact Hall class of
$C_{S_2}$ to \eqref{eq:A4-quantum-variable}.

The earlier exact Hall relation
\[
 [Z_4][U_2]-[U_2][Z_4]=[C_{S_2}]
\]
constructs the differential object.  The new quotient then handles products
of $[C_{S_2}]$ with every incompatible rigid class.  Their non-split and
non-rigid middle terms are straightened by the appropriate heptagon
quadrilateral relations and acquire unique expansions in the $42$ cluster
chambers.

\subsection{A representative crossing}

Choose four cyclically ordered vertices $0<1<3<5$ and the crossing diagonals
\[
 \gamma=(0,3),
 \qquad
 \delta=(1,5).
\]
Their two smoothings are
\[
 (0,1)\cup(3,5),
 \qquad
 (0,5)\cup(1,3).
\]
The first contains the boundary edge $(0,1)$ and hence a frozen unit.  In the
standard normalization the quotient relation is
\begin{equation}
 \overline\eps_{(0,3)}\,
 \overline\eps_{(1,5)}
 =\omega\,
   \overline\eps_{(0,1)\cup(3,5)}
  +\omega^{-1}
   \overline\eps_{(0,5)\cup(1,3)}.
 \label{eq:A4-sample-skein}
\end{equation}
Before quotienting, the left-hand side is an actual exact Hall sum containing
the split crossing object and the non-split middle terms in the selected Hall
direction.  Solving \eqref{eq:A4-sample-skein} for the split object gives its
Hall-corrected theta expansion.

\subsection{Deterministic finite certificate}

The program \texttt{verify\_A4\_hall\_skein.py} works with exact polygon
combinatorics at the classical specialization $\omega=1$.  It verifies:
\begin{center}
\begin{tabular}{lr}
\toprule
quantity & value\\
\midrule
polygon diagonals & \AIVDiagonals\\
quadrilateral crossing pairs & \AIVCrossingPairs\\
triangulations & \AIVTriangulations\\
ordered triple associativity checks & \AIVTripleChecks\\
degree-four multidiagrams checked & \AIVDegreeFourCurves\\
first-crossing confluence comparisons & \AIVConfluenceChecks\\
\bottomrule
\end{tabular}
\end{center}

The triple count is
\[
 14^3=\AIVTripleChecks.
\]
The degree-four run traverses all multisets of four diagonals.  For every
multidiagram with more than one possible first crossing, the program computes
the complete integer-coefficient Ptolemy expansion for every choice and checks
exact equality.

A separate cross-rank program covers cluster types $A_2$ through
$A_{\TypeAMaxRank}$.  Across those \TypeAPolygons\ polygon models it records
\begin{center}
\begin{tabular}{lr}
\toprule
aggregate & value\\
\midrule
diagonals & \TypeATotalDiagonals\\
crossing pairs & \TypeATotalCrossings\\
triangulations & \TypeATotalTriangulations\\
first-choice confluence comparisons & \TypeATotalConfluence\\
\bottomrule
\end{tabular}
\end{center}
using exact integer arithmetic.

\begin{remark}[What the computation certifies]
The code checks the classical polygon state-sum, termination, local confluence
and associativity conventions.  It deliberately does not attach
$\omega^{\pm1}$ to bare unoriented chords: the quantum relation depends on
endpoint elevations/states and, for geometric coefficients, wall data.  The
quantum confluence used in the theorem is supplied by the stated/walled skein
basis theorems.  The code also does not compute every finite-field Hall number
of \eqref{eq:A4-algebra}; those coefficients belong to the exact Hall
expansion and enter the recursive correction
\eqref{eq:Hall-corrected-straightening}.  The proof that their non-split terms
are lower is \cref{thm:extension-smoothing}, not the finite program.
\end{remark}
